\section{实验原理}
\subsection{one-hot法}
\paragraph{基本原理}

One-hot（独热编码）法的核心思想是：为每一个合法的输入输出差分对引入一个二进制选择变量，通过约束所有选择变量之和为 $1$，保证恰好选中一个合法差分对，再利用坐标等式将选中的差分对映射到原始变量，从而实现对合法差分传播集合的精确刻画。

\paragraph{数学建模}

对于一个 $4$ bit 输入、$4$ bit 输出的 S 盒，根据差分分布表（DDT），定义所有合法差分对构成的集合为
\[
V=\left\{
(\Delta_{\mathrm{in}},\Delta_{\mathrm{out}})
\;\middle|\;
\operatorname{DDT}[\Delta_{\mathrm{in}}][\Delta_{\mathrm{out}}]>0
\right\}.
\]

将每个合法差分对展开为 $8$ bit 的二进制向量，记
\[
V=\left\{
v^{(1)},v^{(2)},\ldots,v^{(m)}
\right\}
\subseteq \{0,1\}^{8},
\]
其中 $m=|V|$ 为合法差分对的数量，且
\[
v^{(i)}=
\left(v_1^{(i)},v_2^{(i)},\ldots,v_8^{(i)}\right).
\]

对于集合 $V$ 中的每个合法点 $v^{(i)}$，引入一个二进制选择变量
\[
\lambda_i\in\{0,1\},
\qquad i=1,2,\ldots,m.
\]

首先，要求所有选择变量之和为 $1$：
\[
\boxed{
\sum_{i=1}^{m}\lambda_i=1
}
\]
由于 $\lambda_i$ 均为二进制变量，因此上述约束保证有且仅有一个 $\lambda_i=1$，即恰好选中一个合法差分对。

其次，定义原始差分变量
\[
x=(x_1,x_2,\ldots,x_8)\in\{0,1\}^{8},
\]
并通过坐标等式将其与选择变量关联：
\[
\boxed{
x_j=\sum_{i=1}^{m}v_j^{(i)}\lambda_i,
\qquad j=1,2,\ldots,8.
}
\]

由于 $v_j^{(i)}$ 均为已知的二进制常数，上式是关于变量 $x_j$ 和 $\lambda_i$ 的线性等式。

综上，One-hot 法建立的 MILP 约束模型为
\[
\boxed{
\begin{cases}
\displaystyle\sum_{i=1}^{m}\lambda_i=1,\\[6pt]
\displaystyle x_j=\sum_{i=1}^{m}v_j^{(i)}\lambda_i,
& j=1,2,\ldots,8,\\[6pt]
\lambda_i\in\{0,1\},
& i=1,2,\ldots,m,\\[4pt]
x_j\in\{0,1\},
& j=1,2,\ldots,8.
\end{cases}
}
\]

\paragraph{正确性分析}

由于所有选择变量均为二进制变量，且满足
\[
\sum_{i=1}^{m}\lambda_i=1,
\]
因此必然存在唯一的 $k\in\{1,\ldots,m\}$，使得
\[
\lambda_k=1,\qquad
\lambda_i=0\quad(i\ne k).
\]

代入坐标等式，可得
\[
x_j
=\sum_{i=1}^{m}v_j^{(i)}\lambda_i
=v_j^{(k)},
\qquad j=1,\ldots,8.
\]

因此
\[
x=v^{(k)}\in V,
\]
说明模型不会产生非法差分对。

反之，对于任意合法点 $v^{(k)}\in V$，令
\[
\lambda_k=1,\qquad
\lambda_i=0\quad(i\ne k),
\]
即可满足全部约束，并得到 $x=v^{(k)}$。

因此，该模型在原始变量 $x$ 上的全部可行解恰好构成集合 $V$，即
\[
\boxed{
\left\{
x\in\{0,1\}^{8}
\;\middle|\;
\exists\,\lambda\in\{0,1\}^{m}
\text{ 满足上述约束}
\right\}=V.
}
\]

\subsection{逐点排除法}

\subsection{逻辑化简法}
\subsection{凸包法}
